\chapter{State of the Art}
\label{ch:soa}

This chapter reviews the literature relevant to the two main pillars of
this thesis: the detection of market regimes through statistical and
machine learning methods (\Cref{sec:soa-regimes}), and the extraction of
sentiment from financial text using lexicon-based and deep learning
methods, with particular emphasis on FinBERT and the FNSPID dataset
(\Cref{sec:soa-sentiment}). The chapter closes with a discussion of how
these two strands of research have (rarely) been combined, and identifies
the specific gap addressed by this thesis (\Cref{sec:soa-gap}).

\section{Market Regime Detection}
\label{sec:soa-regimes}

\subsection{Markov-Switching Models and Hidden Markov Models}
\label{sec:soa-hmm}

The idea that the parameters governing the evolution of a financial or
macroeconomic time series can switch between a finite number of
``regimes'' was formalized by Hamilton \cite{hamilton1989}, who proposed a
Markov-switching autoregressive model in which an unobserved discrete
state variable $S_t \in \{1, \dots, K\}$ follows a first-order Markov
chain, and the conditional distribution of the observed series at time
$t$ depends on $S_t$. This framework naturally accommodates the empirical
observation that asset returns exhibit periods of low volatility and
positive drift (\emph{bull} regimes) interspersed with periods of high
volatility and negative drift (\emph{bear} or \emph{crisis} regimes).

A \textbf{Hidden Markov Model (HMM)} is a special case of this framework
in which the observed variables are conditionally independent given the
hidden state, and the state sequence evolves according to a transition
probability matrix $A \in \mathbb{R}^{K \times K}$, with
$A_{ij} = P(S_{t+1}=j \mid S_t = i)$. For continuous observations, each
state $k$ is typically associated with a multivariate Gaussian (or
Gaussian-mixture) emission distribution
$\mathcal{N}(\boldsymbol{\mu}_k, \Sigma_k)$. The classical algorithms for
inference and learning in HMMs --- the forward-backward algorithm for
computing state posteriors, the Viterbi algorithm for the most likely
state sequence, and the Baum-Welch (EM) algorithm for parameter
estimation --- were popularized in the speech-recognition literature by
Rabiner's tutorial \cite{rabiner1989hmm} and have since become a standard
tool in quantitative finance.

In financial applications, Gaussian HMMs have been used extensively for
regime-based asset allocation. Guidolin \cite{guidolin2011} provides a
comprehensive survey of Markov-switching models in empirical finance,
covering applications to asset allocation, derivatives pricing and risk
management. More recently, Nystrup et al. \cite{nystrup2020} demonstrate
that regime-switching models, including HMMs with time-varying transition
probabilities, can be used to dynamically adjust portfolio risk exposure,
improving risk-adjusted returns relative to static allocations. A
recurring theme in this literature is that HMMs naturally model the
\emph{persistence} of regimes --- once the chain enters a given state, it
tends to remain there for several periods, governed by the diagonal
entries of $A$ --- which is a desirable property for regimes that are
expected to last from weeks to months.

However, HMMs also have well-documented limitations. First, the number of
states $K$ must be chosen a priori (or selected via information criteria
such as AIC/BIC), and increasing $K$ rapidly increases the number of free
parameters, particularly for full covariance matrices, which scale as
$O(d^2)$ in the dimensionality $d$ of the observation vector. Second,
because emission distributions are typically assumed Gaussian, HMMs can
struggle to capture regimes characterized by extreme, fat-tailed moves
unless the covariance structure is allowed to be sufficiently flexible.
Third, and most relevant to this thesis, the \emph{smoothness} imposed by
the transition matrix --- which is precisely what gives HMMs their
temporal coherence --- can also cause the model to assign similar,
``averaged'' emission parameters to states that, from an economic
standpoint, would ideally be much more sharply differentiated. This
trade-off between temporal coherence and economic differentiation is
revisited empirically in \Cref{ch:results} of this thesis.

\subsection{Clustering-Based Approaches: K-Means}
\label{sec:soa-kmeans}

An alternative to explicitly modeling temporal transition dynamics is to
treat regime detection as an \emph{unsupervised clustering} problem:
given a set of $N$ feature vectors $\mathbf{x}_t \in \mathbb{R}^d$ (one
per trading day, typically composed of returns, volatility measures,
momentum indicators, etc.), partition them into $K$ groups such that
observations within a group are more similar to each other than to
observations in other groups, \emph{without} explicitly modeling the
sequence of transitions between groups.

The \textbf{K-Means} algorithm \cite{macqueen1967,lloyd1982} is the most
widely used method for this purpose. Given $K$, it seeks to find cluster
assignments $c_t \in \{1, \dots, K\}$ and centroids
$\boldsymbol{\mu}_1, \dots, \boldsymbol{\mu}_K \in \mathbb{R}^d$ that
minimize the within-cluster sum of squares:

\begin{equation}
  \label{eq:soa-kmeans-objective}
  J = \sum_{t=1}^{N} \left\lVert \mathbf{x}_t - \boldsymbol{\mu}_{c_t}
  \right\rVert_2^2 .
\end{equation}

The algorithm alternates between (i) assigning each observation to its
nearest centroid, and (ii) recomputing each centroid as the mean of the
observations assigned to it, until convergence to a local optimum of
\Cref{eq:soa-kmeans-objective}. Because \Cref{eq:soa-kmeans-objective} is
non-convex, the algorithm is typically run multiple times from different
random initializations (e.g., using the \texttt{k-means++}
initialization scheme), and the solution with the lowest objective value
is retained, as implemented in \texttt{scikit-learn}
\cite{pedregosa2011sklearn}.

Compared to HMMs, \textit{K-Means} has two practical advantages for
exploratory regime analysis. First, it does not impose any explicit
temporal structure, which means that the resulting partition reflects
purely the \emph{static} similarity of the feature vectors --- two days
with similar momentum, volatility and RSI values will be assigned to the
same cluster regardless of how far apart they are in time. This can
produce \emph{sharper} clusters in feature space, at the cost of
potentially noisier (less persistent) regime sequences over time. Second,
\textit{K-Means} scales linearly with the number of observations and
features, making it computationally inexpensive even for relatively large
datasets and high-dimensional feature vectors (such as those that include
sector-level sentiment variables, as in this thesis).

The choice of the number of clusters $K$ is typically guided by the
\textbf{elbow method} (plotting $J$ as a function of $K$ and looking for a
point of diminishing returns) and the \textbf{silhouette coefficient},
which for a given observation $i$ is defined as

\begin{equation}
  \label{eq:soa-silhouette}
  s(i) = \frac{b(i) - a(i)}{\max\{a(i), b(i)\}},
\end{equation}

where $a(i)$ is the mean distance between $i$ and the other points in its
own cluster, and $b(i)$ is the smallest mean distance between $i$ and the
points in any other cluster. The silhouette coefficient ranges from $-1$
to $1$, with values close to $1$ indicating well-separated, cohesive
clusters. In financial applications, $K=3$ is a particularly common
choice, as it maps naturally onto the \emph{bull / bear / crisis}
trichotomy that is widely used by practitioners.

\subsection{Limitations of Existing Approaches}
\label{sec:soa-limitations}

Both HMMs and \textit{K-Means} share a common limitation when applied to
market regime detection using only price-derived (technical) features:
the feature vectors are, by construction, \emph{backward-looking}. RSI,
moving-average ratios, momentum and realized volatility are all functions
of \emph{past} prices. As a consequence, regime changes detected by these
models are necessarily \emph{lagged} relative to the underlying economic
or informational shock that triggered them --- the model can only
``see'' a change in regime once it is already reflected in price
dynamics.

This observation motivates the search for \emph{additional} features
that might carry more contemporaneous information about the state of the
market. Sentiment extracted from financial news is a natural candidate:
news about a company, sector or the macroeconomy is, in principle,
available essentially in real time, and may reflect information that has
not yet been fully incorporated into prices. The next section reviews the
literature on sentiment analysis in finance, leading up to the
FinBERT model and the FNSPID dataset used in this thesis.

\section{Sentiment Analysis in Finance}
\label{sec:soa-sentiment}

\subsection{Lexicon-Based Methods}
\label{sec:soa-lexicon}

The earliest approaches to extracting sentiment from financial text
relied on \textbf{word-list (lexicon) based methods}: a document's
sentiment score is computed by counting the occurrences of words that
belong to predefined ``positive'' and ``negative'' word lists, often
normalized by document length. A landmark study by Loughran and McDonald
\cite{loughranmcdonald2011} showed that approximately three-quarters of
the words classified as ``negative'' by the general-purpose Harvard-IV-4
psychosocial dictionary --- a dictionary widely used in earlier finance
studies --- are \emph{not} typically perceived as negative in a financial
context (e.g., words such as ``tax'', ``cost'', ``capital'' or
``liability'' are common in routine financial disclosures and do not, by
themselves, signal bad news). They proposed an alternative,
finance-specific word list, which has since become a standard baseline in
the accounting and finance literature.

While computationally inexpensive and fully interpretable, lexicon-based
methods have well-known limitations: they are insensitive to negation and
context (``did not miss expectations'' may be scored as containing a
negative word), cannot capture sarcasm or implicit sentiment, and require
manual curation of word lists for each domain and, often, each language.

\subsection{Sentiment Indices and Market Outcomes}
\label{sec:soa-sentiment-indices}

A parallel and complementary line of research has focused not on
classifying individual documents, but on constructing \emph{aggregate
sentiment indices} and relating them to subsequent market outcomes.
Tetlock \cite{tetlock2007} analyzed the daily content of a widely-read
\textit{Wall Street Journal} column using the General Inquirer
psychological dictionary, and found that unusually pessimistic media
content predicted downward pressure on market prices, followed by a
reversion to fundamentals, and that high media pessimism was associated
with high subsequent trading volume. Baker and Wurgler
\cite{baker2006} constructed a broader investor sentiment index based on
six market-based proxies (closed-end fund discounts, trading volume,
dividend premia, IPO activity, etc.) and showed that this index has
predictive power for the cross-section of stock returns, particularly for
stocks that are difficult to value or arbitrage.

More recently, the rise of social media has led to a substantial body of
work using sentiment extracted from platforms such as Twitter/X to predict
stock price movements \cite{pagolu2016}, and a broad survey by Xing et al.
\cite{xing2018} categorizes the field of ``natural language based
financial forecasting'', covering data sources (news, social media,
filings), text representations (bag-of-words, embeddings,
Transformer-based), and forecasting targets (price direction, volatility,
risk).

\subsection{Deep Learning and Transformer-Based Sentiment Models}
\label{sec:soa-deep-learning}

The introduction of the \textbf{Transformer} architecture
\cite{vaswani2017attention}, and in particular of pre-trained
bidirectional language models such as \textbf{BERT} (Bidirectional
Encoder Representations from Transformers) \cite{devlin2019bert},
represented a step change in natural language processing. BERT is
pre-trained on large unlabeled text corpora using a masked-language-model
objective, and can subsequently be \emph{fine-tuned} on a relatively small
labeled dataset for a specific downstream task --- such as sentiment
classification --- while retaining the rich contextual representations
learned during pre-training. This approach substantially outperformed
earlier recurrent and convolutional architectures on a wide range of NLP
benchmarks, including sentiment analysis.

In the financial domain, Ding et al. \cite{ding2015} had earlier
demonstrated that deep learning representations of news events (using
structured event tuples and neural embeddings) could improve stock price
movement prediction relative to bag-of-words baselines, foreshadowing the
shift towards learned, contextual representations of financial text that
would later be accelerated by Transformer-based models.

\subsection{FinBERT}
\label{sec:soa-finbert}

\textbf{FinBERT} \cite{finbert2019} is a BERT-based language model further
pre-trained and fine-tuned on financial text (including corporate
filings, earnings call transcripts and analyst reports) for the specific
task of \emph{financial sentiment classification}. Given a piece of
financial text (e.g., a news headline), FinBERT outputs a probability
distribution over three classes --- \emph{positive}, \emph{negative} and
\emph{neutral} --- which can be interpreted as the model's confidence that
the text conveys each type of sentiment from a financial market
perspective.

Empirically, FinBERT and its successors (such as the
\texttt{ProsusAI/finbert} checkpoint used in this thesis, which is
fine-tuned on the Financial PhraseBank dataset of analyst statements) have
been shown to substantially outperform both lexicon-based methods (e.g.,
Loughran-McDonald word counts) and general-purpose sentiment classifiers
on financial text classification benchmarks. Crucially for the present
work, FinBERT can be applied \emph{at scale} to large corpora of news
headlines using GPU-accelerated batched inference, making it feasible to
score millions of headlines spanning more than a decade --- a task that
would be impractical using more computationally expensive
generative large language models, and far less domain-appropriate using
generic lexicons.

In this thesis, FinBERT is used to assign each filtered news headline a
continuous sentiment score $s = p_{\text{pos}} - p_{\text{neg}} \in
[-1, 1]$ and a categorical label given by the arg-max of the three class
probabilities, as described in detail in \Cref{ch:methodology}.

\subsection{The FNSPID Dataset}
\label{sec:soa-fnspid}

\textbf{FNSPID} (Financial News and Stock Price Integration Dataset)
\cite{zhang2023fnspid} is a large-scale, publicly available corpus that
combines several million time-stamped financial news headlines and
articles for US-listed companies with corresponding historical stock
price data. Its scale and temporal coverage make it particularly
well-suited for studies that require aligning news-derived sentiment with
daily price movements over long horizons, such as the present work, which
uses FNSPID news spanning the period from 2009 to 2024.

Compared to smaller, manually curated sentiment datasets (such as the
Financial PhraseBank used to fine-tune FinBERT itself), FNSPID's main
contribution is \emph{scale and temporal alignment}: it provides the raw
material from which a sector-level, time-varying sentiment signal can be
constructed and aligned with technical indicators on a daily basis, which
is a prerequisite for the clustering-based regime detection methodology
proposed in this thesis.

\section{Combining Sentiment and Regime Detection: Related Work and Research Gap}
\label{sec:soa-gap}

A small but growing body of work has begun to combine sentiment-derived
features with regime-switching or clustering models. For example,
Abdollahi et al. \cite{abdollahi2024} and Billah et al.
\cite{billah2024} explore the integration of sentiment signals with
machine learning models for financial forecasting tasks, while Mudarisov
et al. \cite{mudarisov2024} (presented at the ACM International Conference
on AI in Finance, ICAIF) examine related questions at the intersection of
machine learning and financial market analysis. These works provide
encouraging evidence that sentiment-derived features are not redundant
with purely price-based features, but they typically focus on
\emph{supervised} prediction tasks (e.g., forecasting returns or
volatility) rather than on \emph{unsupervised} regime identification, and
rarely provide a systematic \emph{economic} validation of the resulting
groups broken down by sector.

Based on this review, three gaps can be identified that this thesis aims
to address:

\begin{enumerate}
  \item \textbf{Lack of systematic comparison between technical-only and
  technical+sentiment clustering for regime detection.} While sentiment
  indices have been related to market \emph{returns} in a predictive
  (supervised) sense, there is limited work directly comparing
  \textit{unsupervised} regime partitions obtained with and without
  sentiment features, using a controlled experimental design (same
  algorithm, same $K$, same evaluation protocol).
  \item \textbf{Limited sector-level economic validation of detected
  regimes.} Much of the regime-detection literature evaluates models using
  aggregate market-level statistics (e.g., the overall index return and
  volatility per regime). This thesis instead constructs
  equally-weighted sector-level return series for all 11 GICS sectors and
  evaluates each detected regime's economic profile \emph{per sector},
  providing a much more granular picture of what each regime represents
  economically.
  \item \textbf{Lack of direct comparison between K-Means and HMM on the
  same sentiment-augmented feature sets.} Although both \textit{K-Means}
  and HMMs have been applied separately to regime detection, this thesis
  provides a direct, like-for-like comparison of the two methods on
  identical feature sets, with a particular focus on the trade-off between
  the temporal smoothness imposed by the HMM's transition dynamics and the
  sharper, more economically differentiated partitions obtained by
  \textit{K-Means}.
\end{enumerate}

\Cref{ch:methodology} describes in detail how these gaps are addressed
through the design of the data pipeline, the three \textit{K-Means}
models (A, B and C), the Gaussian HMM comparison, and the sector-level
economic validation framework.
